\section{Logic, Sets and Numbers}

\subsection{Logic}

\begin{definition}
    A \textbf{mathematical Statement} is a well-formulated assertion that is strictly either true or false.
\end{definition}

\begin{definition}
    A \textbf{Statement} that depends on one or multiple variables is called a \textbf{predicate} and denoted by $P(x), Q(x), \dots$.
\end{definition}

\begin{definition}
    \textbf{Quantifiers} for predicates:
    \begin{itemize}
        \item $\forall x \in \mathbb{R}: P(x)$ (for all $x$ holds $P(x)$)
        \item $\exists x \in \mathbb{R}: P(x)$ (there exists $x$ s.t. $P(x)$)
        \item $\exists ! x \in \mathbb{R}: P(x)$ (exists exactly one $x$)
        \item $\nexists x \in \mathbb{R}: P(x)$ (exists no $x$ s.t. $P(x)$)
    \end{itemize}
\end{definition}

\begin{important}
    The order of quantifiers matters. $\forall x \exists y$ is not equivalent to $\exists y \forall x$.
\end{important}

\begin{definition}
    Basic \textbf{logical connectives} are defined as:
    $$A \land B \coloneqq A \text{ and } B \text{ are true}$$
    $$A \lor B \coloneqq A \text{ or } B \text{ is true}$$
    $$A \implies B \coloneqq A \text{ implies } B \text{ is true}$$
    $$A \iff B \coloneqq A \text{ is equivalent to } B \text{ is true}$$
\end{definition}

\begin{lemma}
    $$\neg A \lor B \iff A \implies B$$
\end{lemma}

\begin{theorem}
    The following equivalences hold:
    $$\neg (\forall x P(x)) \iff \exists x \neg P(x)$$
    $$\neg (\exists x P(x)) \iff \forall x \neg P(x)$$
    aswell as
    $$\forall x (P(x) \land Q(x)) \iff (\forall x P(x)) \land (\forall x Q(x))$$
    $$\exists x (P(x) \lor Q(x)) \iff (\exists x P(x)) \lor (\exists x Q(x))$$
\end{theorem}

\subsection{Sets}

\begin{definition}
    A \textbf{set} is a collection of elements, that are listed or characterized by a certain property. 
    The empty set := $\emptyset$ and is the only set that contains no elements.
\end{definition}

\begin{definition}
    Let $A$ and $B$ be sets. $A$ is a \textbf{subset} of $B$ denoted by $A \subset B$ iff. 
    $$\forall x (x \in A \implies x \in B)$$
    $A$ is a \textbf{proper subset} of $B$ denoted by $A \subsetneq B$ iff. 
    $$A \subset B \land A \neq B$$
\end{definition}

\begin{definition}
    Let $A$ and $B$ be sets.
    The \textbf{union} of $A$ and $B$ := $A \cup B \coloneqq \{x \mid x \in A \lor x \in B\}$.
\end{definition}

\begin{definition}
    The \textbf{intersection} of $A$ and $B$ := $A \cap B \coloneqq \{x \mid x \in A \land x \in B\}$.
\end{definition}

\begin{definition}
    The \textbf{difference} of $A$ and $B$ := $A \setminus B \coloneqq \{x \mid x \in A \land x \notin B\}$.
\end{definition}

\begin{definition}
    The \textbf{Cartesian product} of $A$ and $B$ := $A \times B \coloneqq \{(a, b) \mid a \in A \land b \in B\}$.
\end{definition}

\begin{remark}
    Note that if $B \subset A$ then the diffrence $A \setminus B$ is called the complement of $B$ in $A$ and := $B^c$.
\end{remark}

\subsection{Numbers}

\begin{remark}
    Some important subsets of the real numbers are 
    $$\mathbb{R}^{+}_{0} / \mathbb{R}^{+}$$ 
    $$\mathbb{R}^{-}_{0} / \mathbb{R}^{-}$$ 
    which define the intervalls 
    $$[0, \infty) / (0, \infty)$$
    $$(-\infty, 0] / (-\infty, 0)$$
\end{remark}

\begin{definition}
    Continuous subsets are called \textbf{intervalls.}
\end{definition}

\begin{definition}
    An intervall is called
    \textbf{open} if it does not contain its endpoints 
    $$I = (a, b) \coloneqq \{x \mid a \lt x \lt b\}$$
    \textbf{closed} if it contains its endpoints 
    $$I = [a, b] \coloneqq \{x \mid a \leq x \leq b\}$$
    \textbf{half-open} if it contains one endpoint and not the other 
    $$I = [a, b) \coloneqq \{x \mid a \leq x \lt b\}$$
\end{definition}

\begin{remark}
    An arbitrary subset of $\mathbb{R}$ is called open if it is a union of open intervalls and closed if its complement is open.
\end{remark}

\begin{definition}
    An \textbf{upper (lower) bound} of a subset $X \subset \mathbb{R}$ is an element $y \in \mathbb{R}$ such that:
    $$\forall x \in X: x \leq (\geq) y$$
    The \textbf{maximum (minimum)} of a subset $X \subset \mathbb{R}$ is the element $x_{0} \in X$ such that:
    $$\forall x \in X: x \leq (\geq) x_{0}$$
    Often denoted as \textbf{max(X)} and \textbf{min(X)}.
\end{definition}

\begin{definition}
    An intervall $I = (a, b)$ is \textbf{bounded} if there exist upper and lower bounds hence 
    $$a > -\infty \quad \text{and} \quad b < \infty$$
\end{definition}

\begin{remark}
    A bounded and closed intervall is often refered to as \textbf{compact}.
\end{remark}

\begin{definition}
    The \textbf{Supremum} $sup(X)$ of a subset $X \subset \mathbb{R}$ is the least upper bound of $X$. and can be characterized by:
    $$\forall x \in X: x \leq sup(X) \land$$
    $$\forall \epsilon \gt 0: \exists x \in X: x \gt sup(X) - \epsilon$$    
\end{definition}

\begin{definition}
    The \textbf{Infimum} $inf(X)$ of a subset $X \subset \mathbb{R}$ is the greatest lower bound of $X$. and can be characterized by:
    $$\forall x \in X: x \geq inf(X) \land$$
    $$\forall \epsilon \gt 0: \exists x \in X: x \lt inf(X) + \epsilon$$
\end{definition}

\begin{theorem}
    The \textbf{maximum and minimum are unique} aslong as they exist.
\end{theorem}

\begin{definition}
    \textbf{Order Axioms}
    reflexivity: 
    $$\forall x \in X: x \leq x$$
    antisymmetry: 
    $$\forall x, y \in X: x \leq y \land y \leq x \implies x = y$$
    transitivity: 
    $$\forall x, y, z \in X: x \leq y \land y \leq z \implies x \leq z$$
    totality: 
    $$\forall x, y \in X: x \leq y \lor y \leq x$$
\end{definition}

\begin{definition}
    A set $X$ is said to be \textbf{(order)complete} iff for every non empty subsets $A, B \subseteq X$
    $$\forall a \in A \forall b \in B: a \leq b$$
    $$\implies \exists c \in X \quad \text{s.t.}$$
    $$\quad \forall a \in A: a \leq c \land$$
    $$\forall b \in B: c \leq b$$
\end{definition}

\begin{example}
    The real numbers $\mathbb{R}$ satisfy axiomatic properties for addition, multiplication, distributivity and order. Furthermore $\mathbb{R}$ is an order complete field.
\end{example}

\begin{example}
    $\mathbb{Q}$ is not order complete. Let $A = \{x \in \mathbb{Q} \mid x^2 \leq 2\}$ and $B = \{x \in \mathbb{Q} \mid x^2 \geq 2\}$. Then $\forall a \in A \forall b \in B: a \leq b$ but there is no $c \in \mathbb{Q}$ such that $c^2 = 2$.
\end{example}

\begin{theorem}
    Every \textbf{non empty, bounded from above (below)} subset $X \subset \mathbb{R}$ has a \textbf{supremum (infimum)}.
\end{theorem}

\begin{theorem}
    The \textbf{Archimedian Principle} says that
    For $x \in \mathbb{R}$ and $y \gt 0$ there exists an element 
    $$n \in \mathbb{N} \quad \text{s.t.} \quad n \cdot y \gt x$$
    For every $\epsilon \gt 0$ there exists an element $n \in \mathbb{N}$ such that
    $$\frac{1}{n} \lt \epsilon$$
\end{theorem}

\begin{theorem}
    The Young Inequality states that for all $\epsilon \gt 0$ and all $x, y \in \mathbb{R}$:
    $$2\lvert xy \rvert \leq \epsilon x^2 + \frac{1}{\epsilon} y^2$$
\end{theorem}

\subsubsection{Vectors and Vector Spaces}

\begin{definition}

    The \textbf{euclidian space} $\mathbb{R}^n$ is the set of all vectors $x = (x_1, x_2, \dots, x_n)$ where $x_i \in \mathbb{R}$.

\end{definition}

\begin{definition}

    \textbf{Addition} and \textbf{scalar multiplication} are defined as:
    \begin{align*}
& x + y \\
&= (x_1 + y_1, x_2 + y_2, \dots, x_n + y_n)
\end{align*}
    \begin{align*}
& \lambda x \\
&= (\lambda x_1, \lambda x_2, \dots, \lambda x_n)
\end{align*}

\end{definition}

\begin{definition}
    A \textbf{vector space} must satisfy the following axioms: \\
    Commutativity: $x + y = y + x$ \\
    Associativity: $x + (y + z) = (x + y) + z$ \\
    Identity: $\exists e \in X \text{ s.t. } x + e = x \quad \forall x \in X$ \\
    Inverse: $\forall x \in X \exists x^{-1} \in X \text{ s.t. } x + x^{-1} = e$ \\
    Distributivity: $x \cdot (y + z) = x \cdot y + x \cdot z$
\end{definition}

\begin{definition}
    For $x, y \in \mathbb{R}^n$, the \textbf{standard scalar product} is :=:
    $$x \cdot y = \langle x, y \rangle \coloneqq \sum_{i=1}^n x_i y_i$$
    The \textbf{euclidean norm} is:
    $$||x|| \coloneqq \sqrt{\sum_{i=1}^n x_i^2}$$
    The \textbf{euclidean distance} between $x$ and $y$ is:
    $$d(x, y) \coloneqq \sqrt{\sum_{i=1}^n (x_i - y_i)^2}$$

\end{definition}

\begin{theorem}
    The triangle- and the cauch-schwarz inequality hold for all $x, y, z \in \mathbb{R}^n$:
    $$||x - z|| \leq ||x - y|| + ||y - z||$$
    $$|\langle x, y \rangle| \leq ||x|| ||y||$$
\end{theorem}

\subsubsection{Complex Numbers}

\begin{definition}
    The \textbf{complex numbers} $\mathbb{C}$ are the set of numbers
    $$ \mathbb{C} = \{a + ib \mid a, b \in \mathbb{R}\}$$
    where $a$ is the \textbf{real part} $\text{Re}(z)$, and $b$ is the \textbf{imaginary part} $\text{Im}(z)$.
\end{definition}

\begin{definition}
    Arithmetic of complex numbers $z_1 = x_1 + i y_1$ and $z_2 = x_2 + i y_2$:
    Addition: 
    $$z_1 + z_2 = (x_1 + x_2) + i(y_1 + y_2)$$
    Multiplication: 
    $$z_1 \cdot z_2 = (x_1 x_2 - y_1 y_2) + i(x_1 y_2 + y_1 x_2)$$
\end{definition}

\begin{definition}
    The \textbf{complex conjugate} of $z = x + iy$ is $\bar{z} = x - iy$.
\end{definition}

\begin{important}
    Division is calculated by expanding with the conjugate of the denominator:
    $$\frac{z}{w} = \frac{z \cdot \bar{w}}{w \cdot \bar{w}} = \frac{z \cdot \bar{w}}{|w|^2}$$
\end{important}

\begin{lemma}
    Some usfull properties of the conjugate:
    $$z \cdot \bar{z} = |z|^2$$
    $$z + \bar{z} = 2\text{Re}(z)$$
    $$z - \bar{z} = 2i\text{Im}(z)$$
    $$\overline{z+w} = \bar{z} + \bar{w}$$
    $$\overline{z \cdot w} = \bar{z} \cdot \bar{w}$$
\end{lemma}

\begin{definition}
    Geometrically, complex numbers can be represented in the Gaussian plane. The \textbf{absolute value} of $z = a + ib$ is its \textbf{distance from the origin}:
    $$|z| = \sqrt{a^2 + b^2}$$
    The distance between $z_1$ and $z_2$ is $d = |z_1 - z_2|$. 
\end{definition}

\begin{theorem}
    The Euler's formula states that for $t \in \mathbb{R}$:
    $$e^{it} = \cos(t) + i\sin(t)$$
    This allows writing trigonometric functions using complex exponentials:
    $$\cos(t) = \frac{e^{it} + e^{-it}}{2}$$
    $$\sin(t) = \frac{e^{it} - e^{-it}}{2i}$$
\end{theorem}

\begin{remark}
    The \textbf{Intuition} behind eulers formula is we can use Taylor series to show that:
    $$e^x = 1 + x + \frac{x^2}{2!} + \frac{x^3}{3!} + \dots = \sum_{n=0}^\infty \frac{x^n}{n!}$$
    Using $x=it$ we get:
    $$e^{it} = \sum_{n=0}^\infty \frac{(it)^n}{n!} = \dots = \cos(t) + i\sin(t)$$
\end{remark}

\begin{definition}
    For $z \in \mathbb{C}$, the polar-coordinates are defined as:
    $$z = r(\cos(t) + i\sin(t)) = r e^{it}$$
    where $r = |z|$ is the absolute value and $t = \arg(z)$ is the argument.
\end{definition}

\begin{remark}
    To convert from cartesian to polar coordinates:
    $$r = \sqrt{a^2 + b^2}$$
    $$t = \arctan\left(\frac{b}{a}\right)$$
    Note that the arctan function only returns values in $(-\frac{\pi}{2}, \frac{\pi}{2})$, so we have to adjust the angle based on the quadrant of the complex number.
\end{remark}

\begin{example}
    Multiplication and division in polar-coordinates is much easier than in cartesian coordinates:
    $$z_1 \cdot z_2 = r_1 e^{it_1} \cdot r_2 e^{it_2} = r_1 r_2 e^{i(t_1 + t_2)}$$
    $$\frac{z_1}{z_2} = \frac{r_1 e^{it_1}}{r_2 e^{it_2}} = \frac{r_1}{r_2} e^{i(t_1 - t_2)}$$
\end{example}

\begin{example}
    So is exponentiation or taking roots:
    $$z^n = (r e^{it})^n = r^n e^{int}$$
    $$\sqrt[n]{z} = \sqrt[n]{r} e^{i\frac{t+2\pi k}{n}}, \quad k \in \{0, 1, \dots, n-1\}$$
\end{example}

\begin{definition}
    The $n$-th roots of unity are the solutions to $z^n = 1$, which are given by:
    $$z_k = e^{i\frac{2\pi k}{n}}, \quad k \in \{0, 1, \dots, n-1\}$$
\end{definition}

\begin{theorem}
    The sum of the roots of unit is always zero.
    $$\sum_{k=0}^{n-1} e^{i\frac{2\pi k}{n}} = 0$$
\end{theorem}

\begin{theorem}
    The \textbf{fundamental theorem of algebra} states that every polynomial of degree $n \geq 1$ has exactly $n$ roots in $\mathbb{C}$ (counting multiplicity). These roots appear in konplex-conjugate pairs.
\end{theorem}
\subsection{Landau Notation (Asymptotics)}
\begin{definition}
    Let $f, g$ be functions.
    \begin{itemize}
        \item $f(x) = O(g(x))$ as $x \to x_0$ iff $\exists C, \delta > 0$ s.t. $|f(x)| \le C|g(x)|$ for $|x-x_0| < \delta$.
        \item $f(x) = o(g(x))$ as $x \to x_0$ iff $\lim_{x \to x_0} \frac{f(x)}{g(x)} = 0$.
        \item $f(x) \sim g(x)$ as $x \to x_0$ iff $\lim_{x \to x_0} \frac{f(x)}{g(x)} = 1$.
    \end{itemize}
\end{definition}

\begin{remark}
    \textbf{Useful for Taylor:} $f(x) = P_n(x) + O(|x-x_0|^{n+1})$.
\end{remark}

\subsection{Exam Strategy: Sets and Topology}
\begin{itemize}
    \item \textbf{Open:} Every point is an interior point (ball $B_\epsilon(x) \subset A$). Check: Is the boundary excluded?
    \item \textbf{Closed:} Complement is open OR contains all limit points. Check: Is the boundary included?
    \item \textbf{Bounded:} $A \subset B_R(0)$ for some $R$.
    \item \textbf{Compact:} In $\mathbb{R}^n$, compact $\iff$ closed and bounded.
    \item \textbf{Supremum/Infimum Tips:} 
        \begin{itemize}
            \item For $A = \{a_n : n \in \mathbb{N}\}$, check if $a_n$ is monotonic. If $a_n \nearrow L$, then $\sup A = L$, $\inf A = a_1$.
            \item If $a_n$ oscillates, check even/odd indices separately.
        \end{itemize}
\end{itemize}

\subsection{Exam Strategy: Complex Equations}
\begin{itemize}
    \item \textbf{Circles:} $|z - z_0| = R$ is a circle of radius $R$ around $z_0$.
    \item \textbf{Lines:} $|z - a| = |z - b|$ is the perpendicular bisector of $a$ and $b$.
    \item \textbf{Polynomials:} If $P(z)=0$ and coefficients are real, then $z \in \mathbb{C}$ is a root $\implies \bar{z}$ is also a root.
    \item \textbf{Roots of Unity Sum:} $\sum_{k=0}^{n-1} z_k = 0$ for $n > 1$.
\end{itemize}
